\documentclass[10pt,a4paper]{amsart}
\usepackage{amsmath,amssymb,amsthm,graphicx}
\usepackage{verbatim}
\newenvironment{reference}{\comment}{\endcomment}
%\newenvironment{reference}{}{}
\newenvironment{slogan}{\comment}{\endcomment}
\newenvironment{history}{\comment}{\endcomment}

% For commutative diagrams we use Xy-pic
\usepackage[all]{xy}

% We use 2cell for 2-commutative diagrams.
\xyoption{2cell}
% Theorem environments.
%
\theoremstyle{plain}
\newtheorem{theorem}[subsection]{Theorem}
\newtheorem{proposition}[subsection]{Proposition}
\newtheorem{lemma}[subsection]{Lemma}

\theoremstyle{definition}
\newtheorem{definition}[subsection]{Definition}
\newtheorem{example}[subsection]{Example}
\newtheorem{exercise}[subsection]{Exercise}
\newtheorem{situation}[subsection]{Situation}

\theoremstyle{remark}
\newtheorem{remark}[subsection]{Remark}
\newtheorem{remarks}[subsection]{Remarks}

\numberwithin{equation}{subsection}

% Macros
%
\def\lim{\mathop{\mathrm{lim}}\nolimits}
\def\colim{\mathop{\mathrm{colim}}\nolimits}
\def\Spec{\mathop{\mathrm{Spec}}}
\def\Hom{\mathop{\mathrm{Hom}}\nolimits}
\def\Ext{\mathop{\mathrm{Ext}}\nolimits}
\def\SheafHom{\mathop{\mathcal{H}\!\mathit{om}}\nolimits}
\def\SheafExt{\mathop{\mathcal{E}\!\mathit{xt}}\nolimits}
\def\Sch{\mathit{Sch}}
\def\Mor{\mathop{\mathrm{Mor}}\nolimits}
\def\Ob{\mathop{\mathrm{Ob}}\nolimits}
\def\Sh{\mathop{\mathit{Sh}}\nolimits}
\def\NL{\mathop{N\!L}\nolimits}
\def\CH{\mathop{\mathrm{CH}}\nolimits}
\def\proetale{{pro\text{-}\acute{e}tale}}
\def\etale{{\acute{e}tale}}
\def\QCoh{\mathit{QCoh}}
\def\Ker{\mathop{\mathrm{Ker}}}
\def\Im{\mathop{\mathrm{Im}}}
\def\Coker{\mathop{\mathrm{Coker}}}
\def\Coim{\mathop{\mathrm{Coim}}}

% Boxtimes
%
\DeclareMathSymbol{\boxtimes}{\mathbin}{AMSa}{"02}

%
% Macros for moduli stacks/spaces
%
\def\QCohstack{\mathcal{QC}\!\mathit{oh}}
\def\Cohstack{\mathcal{C}\!\mathit{oh}}
\def\Spacesstack{\mathcal{S}\!\mathit{paces}}
\def\Quotfunctor{\mathrm{Quot}}
\def\Hilbfunctor{\mathrm{Hilb}}
\def\Curvesstack{\mathcal{C}\!\mathit{urves}}
\def\Polarizedstack{\mathcal{P}\!\mathit{olarized}}
\def\Complexesstack{\mathcal{C}\!\mathit{omplexes}}
% \Pic is the operator that assigns to X its picard group, usage \Pic(X)
% \Picardstack_{X/B} denotes the Picard stack of X over B
% \Picardfunctor_{X/B} denotes the Picard functor of X over B
\def\Pic{\mathop{\mathrm{Pic}}\nolimits}
\def\Picardstack{\mathcal{P}\!\mathit{ic}}
\def\Picardfunctor{\mathrm{Pic}}
\def\Deformationcategory{\mathcal{D}\!\mathit{ef}}
\usepackage{cleveref}
\crefname{originalreference}{original reference}{original references}
\Crefname{originalreference}{Original reference}{Original references}
\makeatletter
\newlabel{sites-lemma-directed-colimits-sections}{{0738}{source}{}{sourcepointer-1}{}}
\newlabel{sites-lemma-directed-colimits-sections@cref}{{[originalreference][0][]0738}{[1][0][]source}{}{}{}}
\makeatother
\crefname{originallemma}{original lemma}{original lemmas}
\Crefname{originallemma}{Original lemma}{Original lemmas}
\makeatletter
\newlabel{lemma-limits-colimits-abelian-sheaves}{{03CO}{source}{}{sourcepointer-2}{}}
\newlabel{lemma-limits-colimits-abelian-sheaves@cref}{{[originallemma][0][]03CO}{[1][0][]source}{}{}{}}
\makeatother
\title{Sections over a quasi-compact object commute with coproducts and direct sums}
\author{The Stacks Project Authors}
\date{stacks-project:sites-modules.tex:local-snapshot-20261012; GNU FDL 1.2 or later}
\begin{document}
\maketitle
\noindent\textit{Editorial note.} All delivered mathematical text is the author's own, copied verbatim from the pinned Stacks Project chapter file sites-modules.tex; no statement or proof step was rewritten, reordered or removed.\par\medskip
\noindent\textit{Editorial note.} The counted span cites two external items (the directed-colimits lemma and the description of filtered colimits of abelian sheaves), recorded in assembly.external\_references with canonical tags as navigation pointers.\par\medskip
\noindent\textit{Allowed standard prerequisites (editorial).}\begin{itemize}
\item Sections over a quasi-compact object commute with filtered colimits with injective transition maps. Stacks Project, Sites, Lemma sites-lemma-directed-colimits-sections (tag 00YQ). Supplies the commutation result to which both statements are reduced.
\end{itemize}\medskip
\ifdefined\hypertarget\hypertarget{sourcepointer-1}{}\fi
\noindent\textit{Source reference.} Original reference 0738: sources/stacks-chow/tags.txt, source lines 9052--9052. Cross-chapter standard prerequisite cited by the author inside the counted span; it is used as a pointer and its own proof is not part of this item.\par\medskip
\ifdefined\hypertarget\hypertarget{sourcepointer-2}{}\fi
\noindent\textit{Source reference.} Original lemma 03CO: sources/stacks-chow/tags.txt, source lines 4391--4391. Cross-chapter standard prerequisite cited by the author inside the counted span; it is used as a pointer and its own proof is not part of this item.\par\medskip
\begin{lemma}
\label{lemma-sections-over-quasi-compact}
Let $\mathcal{C}$ be a site. Let $W$ be a quasi-compact
object of $\mathcal{C}$.
\begin{enumerate}
\item The functor $\Sh(\mathcal{C}) \to \textit{Sets}$,
$\mathcal{F} \mapsto \mathcal{F}(W)$ commutes with coproducts.
\item Let $\mathcal{O}$ be a sheaf of rings on $\mathcal{C}$. The functor
$\textit{Mod}(\mathcal{O}) \to \textit{Ab}$,
$\mathcal{F} \mapsto \mathcal{F}(W)$
commutes with direct sums.
\end{enumerate}
\end{lemma}

\begin{proof}
Proof of (1). Taking sections over $W$ commutes with filtered colimits
with injective transition maps by
Sites, Lemma \ref{sites-lemma-directed-colimits-sections}.
If $\mathcal{F}_i$ is a family of sheaves of sets
indexed by a set $I$. Then $\coprod \mathcal{F}_i$ is the filtered
colimit over the partially ordered set of finite subsets
$E \subset I$ of the coproducts
$\mathcal{F}_E = \coprod_{i \in E} \mathcal{F}_i$.
Since the transition maps are injective we conclude.

\medskip\noindent
Proof of (2).
Let $\mathcal{F}_i$ be a family of sheaves of $\mathcal{O}$-modules
indexed by a set $I$. Then $\bigoplus \mathcal{F}_i$ is the filtered
colimit over the partially ordered set of finite subsets
$E \subset I$ of the direct sums
$\mathcal{F}_E = \bigoplus_{i \in E} \mathcal{F}_i$.
A filtered colimit of abelian sheaves can be computed in the
category of sheaves of sets. Moreover, for $E \subset E'$ the transition map
$\mathcal{F}_E \to \mathcal{F}_{E'}$ is injective (as sheafification
is exact and the injectivity is clear on underlying presheaves).
Hence it suffices to show the result for a finite index set by
Sites, Lemma \ref{sites-lemma-directed-colimits-sections}.
The finite case is dealt with in
Lemma \ref{lemma-limits-colimits-abelian-sheaves}
(it holds over any object of $\mathcal{C}$).
\end{proof}

\end{document}
